\documentclass[11pt,letterpaper]{article}
\usepackage{amsmath,amssymb,amsthm}
\begin{document}
\section{Propositional logic}
We write $\neg$, $\wedge$, $\vee$, $\to$ for negation, conjunction, disjunction and implication. A formula $\varphi$ is a \emph{tautology} if it is true under every valuation.

\begin{quote}
A proof is whatever convinces a reasonable reader; a formal proof is whatever convinces a machine.
\end{quote}

\subsection{An example}
The formula $(p\to q)\to(\neg q\to\neg p)$ is a tautology, as the table shows.
\begin{center}
\begin{tabular}{cc|c|c|c}
$p$ & $q$ & $p\to q$ & $\neg q\to\neg p$ & whole\\ \hline
T & T & T & T & T\\
T & F & F & F & T\\
F & T & T & T & T\\
F & F & T & T & T
\end{tabular}
\end{center}

\subsection{Natural deduction}
\begin{enumerate}
\item Assume $p\to q$ and $\neg q$.
\item Assume $p$; then $q$ by modus ponens.
\item This contradicts $\neg q$, so $\neg p$.
\end{enumerate}
\smallskip
The verbatim sketch of the same derivation:
\begin{verbatim}
  1. p -> q      assumption
  2. ~q          assumption
  3.   p         assumption
  4.   q         MP 1,3
  5.   false     2,4
  6. ~p          RAA 3-5
\end{verbatim}
Here \verb|RAA| stands for reductio ad absurdum.
\end{document}
